\section{Desperate Measures}

In order to make sensible statistical statements about the prevalence of eternal inflation, we must identify the set $\Omega$ of possible outcomes under consideration, and then establish a system for assigning probabilities to those various outcomes.  
In our context, we take $\Omega$ to be the set of viable inflation scenarios -- pairings of a potential function and initial conditions.  

The probability {\it measure} is a map $\mu : 2^\Omega \to \mathbb R$ from possible sets of draws from the ``hat'' $\Omega$ to the interval $[0,1]$.  The measure assigns a probability to each set of outcomes, such that $\mu(\bigwedge_i x_i) = \sum_i \mu(x_i)$ for disjoint sets $x_i \in 2^\Omega$,
% $\mu(x) \leq \mu(x \bigcup y) \,\forall\, x,y \in 2^\Omega$, 
with $\mu(\emptyset)=0$ and $\mu(\Omega) = 1$.
Unlike the measures on the volume of Euclidean space ($\dx x \wedge \dx y \wedge \dx z$) or on rolls of a weighted die, the measure in our case is ambiguous; there is no physically preferred way of assigning probabilities to the space of inflation scenarios.

Can we find a measure in which eternal inflation is not generic, which is reasonably simple and not cheating?

What about the No-Boundary Measure or a measure that exponentially favors low energies.

It is likely not reasonable to expect that the measures on the potential $V(\phi)$ and the initial conditions $\mathbf p$ are separable, such that their joint measure would take the form $\mu(V,\mathbf p) = \mu(V) \mu(\mathbf p)$.  (A counterexample might be found within the subspace consisting of quadratic potentials, in which arbitrarily large field values may be disfavored on physical grounds.)





\subsection{Measures on Potentials}

If the inflaton is to be treated as a fundamental quantum field, then the classical potential to which the expectation value of the inflaton is approximately subject will the expectation value of an operator-valued function of the field operator: $V(\phi) \approx \bra{\Psi} V(\hat \phi) \ket{\Psi}$.  If the field is allowed to take values comparable to the Planck mass, then we may construct a potential with any number of terms containing even or odd powers of the field operator, so long as the potential is bounded from below.  
% If such a potential possesses broken symmetry, then expanding the potential around the nonzero vacuum expectation value will likely yield terms that are odd powers of the field operator.  
If the field is coupled to additional fields with vacuum expectation values and that vary slowly in time, then such interaction terms would be even-powered.


While the potential function may contain odd terms in the series expansion, it must be bounded from below.  Even if we restrict the function space of potentials to the square integrable functions $L^2$.


\begin{itemize}
    
    \item Weigh potentials by their complexity -- by the number of interactions with unmodeled free fields that would be necessary to produce the potential.
    \item Weigh potentials by their degree of departure from plausible potentials for elementary fields in quantum field theory.

\end{itemize}

If we assume that the inflaton is a quantum scalar field expanded about some minimum in its potential, what assumptions can we then make about the form of that potential?

\paragraph{Reconstructing the Potential}

\subsection{Measures on Initial Conditions}

    The initial conditions in our models are the values of the homogeneous scalar field $\phi$ and its time derivative.  A simple example of a measure over initial field values is $\mu(\Delta) = \int_\Delta \dx \phi$, where $\Delta$ is some finite (or infinitesimal) interval in field space.  This measure assigns equal probability to all field values.
    
    Tegmark \cite{tegmark08} adopted two measures on the space of initial conditions:
    \begin{description}
    
        \item[Measure A] {\it Assign equal marginal probability to the field values that maximize the potential function, and zero probability to all other field values.} %, unit marginal probability to $\dot\phi = 0$.
        
        In this case, we are sampling half-basins of the potential, accounting for any inflation that occurs between a randomly selected local maximum and a neighboring minimum.  We allow the field to approach the local maximum as $\tau \to -\infty$.
        
        \begin{description}
        \item[Pros] $\dot\phi_0$ given by slow roll equation of motion.
        \item[Cons] There may be inflation higher on the potential, or the inflation epoch into which the field is dropped may have already generated enough e-folds of expansion to satisfy the bound on $N$.
        \end{description}
        
        \item[Measure B] {\it Assign equal marginal probability to all field values.  Discard scenarios in which slow roll conditions are not met at $\phi_0$. % unit marginal probability to $\dot\phi = 0$.
        }
        
        This is equivalent to dropping the field at a random point and starting the clock wherever it lands.  Since in the end we intend to condition our probabilities on the presence of inflation, we lose nothing by discarding scenarios in which the initial field value does not fall within an inflating region of the potential.
        
        \begin{description}
        \item[Pros] $\dot\phi_0$ given by slow roll equation of motion.
        \item[Cons] There may be inflation higher on the potential, or the inflation epoch into which the field is dropped may have already generated enough e-folds of expansion to satisfy the bound on $N$.
        \end{description}
        
    \end{description}
    
    These measures were chosen in part because they allow us to set $\dot\phi$ to zero while maintaining some peace of mind.
    
    Adopting Measure B would lead to a unit probability of encourering stochastic eternal inflation, since the criterion \eqref{seic} is always satisfied when $V_{,\phi} \to 0$.

